\documentclass[11pt]{article}
\usepackage[a4paper,margin=1in]{geometry}
\usepackage{fontspec}
\usepackage[boldfont=false]{xeCJK}
\setCJKmainfont{FandolSong-Regular.otf}
\usepackage{amsmath}
\usepackage{amssymb}
\usepackage{array}
\usepackage{booktabs}
\usepackage{graphicx}
\usepackage{adjustbox}
\usepackage{enumitem}
\setlist[itemize]{topsep=2pt,partopsep=0pt,parsep=0pt,itemsep=2pt,leftmargin=1.4em}
\setlist[enumerate]{topsep=2pt,partopsep=0pt,parsep=0pt,itemsep=2pt,leftmargin=1.6em}
\usepackage{parskip}
\usepackage{fvextra}
\fvset{breaklines=true,breakanywhere=true,fontsize=\small}
\setlength{\parindent}{0pt}

\begin{document}

\begin{center}
  {\LARGE\bfseries Boolean logic}\\[0.35em]
  {\large Igcse Computer Science}
\end{center}
\vspace{0.6em}

\section*{What is Boolean logic?}

\textbf{Boolean logic} 布尔逻辑 works with values that are either \textbf{true} or \textbf{false}. In electronics these are shown as \textbf{1} (true) and \textbf{0} (false). A \textbf{logic gate} 逻辑门 takes one or more of these inputs and gives one output, following a fixed rule.

A \textbf{truth table} 真值表 lists every possible set of inputs and the output for each. You build it by writing out all the input combinations.

\section*{The six logic gates}

You must know six gates. \texttt{NOT} has one input; all the others have two inputs (A and B).

\subsection*{NOT gate}

The \textbf{NOT gate} 非门 reverses the input. Output is 1 when the input is 0.

\begin{center}
\adjustbox{max width=\linewidth}{%
\begin{tabular}{ll}
\toprule
\textbf{A} & \textbf{Output} \\
\midrule
0 & 1 \\
1 & 0 \\
\bottomrule
\end{tabular}}
\end{center}

\subsection*{AND gate}

The \textbf{AND gate} 与门 gives output 1 \textbf{only when both} inputs are 1.

\begin{center}
\adjustbox{max width=\linewidth}{%
\begin{tabular}{lll}
\toprule
\textbf{A} & \textbf{B} & \textbf{Output} \\
\midrule
0 & 0 & 0 \\
0 & 1 & 0 \\
1 & 0 & 0 \\
1 & 1 & 1 \\
\bottomrule
\end{tabular}}
\end{center}

\subsection*{OR gate}

The \textbf{OR gate} 或门 gives output 1 when \textbf{at least one} input is 1.

\begin{center}
\adjustbox{max width=\linewidth}{%
\begin{tabular}{lll}
\toprule
\textbf{A} & \textbf{B} & \textbf{Output} \\
\midrule
0 & 0 & 0 \\
0 & 1 & 1 \\
1 & 0 & 1 \\
1 & 1 & 1 \\
\bottomrule
\end{tabular}}
\end{center}

\subsection*{NAND gate}

The \textbf{NAND gate} 与非门 is AND followed by NOT. The output is the \textbf{opposite} of AND.

\begin{center}
\adjustbox{max width=\linewidth}{%
\begin{tabular}{lll}
\toprule
\textbf{A} & \textbf{B} & \textbf{Output} \\
\midrule
0 & 0 & 1 \\
0 & 1 & 1 \\
1 & 0 & 1 \\
1 & 1 & 0 \\
\bottomrule
\end{tabular}}
\end{center}

\subsection*{NOR gate}

The \textbf{NOR gate} 或非门 is OR followed by NOT. The output is the \textbf{opposite} of OR.

\begin{center}
\adjustbox{max width=\linewidth}{%
\begin{tabular}{lll}
\toprule
\textbf{A} & \textbf{B} & \textbf{Output} \\
\midrule
0 & 0 & 1 \\
0 & 1 & 0 \\
1 & 0 & 0 \\
1 & 1 & 0 \\
\bottomrule
\end{tabular}}
\end{center}

\subsection*{XOR gate}

The \textbf{XOR gate} 异或门 (exclusive OR) gives output 1 when the inputs are \textbf{different}.

\begin{center}
\adjustbox{max width=\linewidth}{%
\begin{tabular}{lll}
\toprule
\textbf{A} & \textbf{B} & \textbf{Output} \\
\midrule
0 & 0 & 0 \\
0 & 1 & 1 \\
1 & 0 & 1 \\
1 & 1 & 0 \\
\bottomrule
\end{tabular}}
\end{center}

\section*{Logic expressions}

A \textbf{logic expression} 逻辑表达式 writes a circuit using letters and gate words. The usual way to write the gates:

\begin{center}
\adjustbox{max width=\linewidth}{%
\begin{tabular}{ll}
\toprule
\textbf{Gate} & \textbf{In words} \\
\midrule
NOT A & NOT A \\
A AND B & A AND B \\
A OR B & A OR B \\
\bottomrule
\end{tabular}}
\end{center}

For example, the expression \texttt{(A AND B) OR (NOT C)} means: do A AND B, do NOT C, then OR the two results together.

\section*{Logic circuits}

A \textbf{logic circuit} 逻辑电路 joins gates together to carry out a task. The output of one gate can become the input of another. At IGCSE a circuit has up to \textbf{three inputs} and \textbf{one output}.

You must be able to move between three forms:

\begin{itemize}
\item a \textbf{problem statement} 问题陈述 (a description in words),

\item a logic expression,

\item a logic circuit,

\item a truth table.

\end{itemize}

\subsection*{From a problem statement to a circuit}

Read the statement and pick out the conditions and the logic words (and, or, not). For example:

\begin{quote}
An alarm (X) sounds when the door is open (A) AND the system is switched on (B).

\end{quote}

This is \texttt{X = A AND B}, so you draw one AND gate with inputs A and B.

\subsection*{Completing a truth table from a circuit or expression}

To fill in a truth table:

\begin{enumerate}
\item Write \textbf{all} the input combinations. For three inputs there are 8 rows (000 up to 111).

\item Work out each gate's output in order, one column at a time.

\item The last column is the final output.

\end{enumerate}

\begin{center}
\adjustbox{max width=\linewidth}{%
\begin{tabular}{lllll}
\toprule
\textbf{A} & \textbf{B} & \textbf{C} & \textbf{A AND B} & \textbf{(A AND B) OR C} \\
\midrule
0 & 0 & 0 & 0 & 0 \\
0 & 0 & 1 & 0 & 1 \\
0 & 1 & 0 & 0 & 0 \\
0 & 1 & 1 & 0 & 1 \\
1 & 0 & 0 & 0 & 0 \\
1 & 0 & 1 & 0 & 1 \\
1 & 1 & 0 & 1 & 1 \\
1 & 1 & 1 & 1 & 1 \\
\bottomrule
\end{tabular}}
\end{center}

Adding a middle "working" column for each gate makes the final output easy to fill in. Always draw the circuit exactly as the statement says, \textbf{without simplifying} it.


\end{document}
